\src{192}
\subsection*{REMARKS}

\noindent 1. For the definitions preceding Theorem 6, it was not necessary for
$X$ to be nonsingular, since it was not essential to work only with finite
resolutions. But if $X$ is singular, one can no longer assert a priori that the
$(\underline{T}_{r}^{s})^{p}(A_1,\ldots;B_1,\ldots)$ are coherent sheaves,
because the complex of sheaves
\[
  \Homsh_{\Osh_X}\bigl(L(A_1)\otimes\cdots,L(B_1)\otimes\cdots\bigr)
\]
will have infinitely many components of any given total degree.

\medskip
\noindent 2. It is readily checked that in the formulas (8.1), one of the
$L(B_i)$ may be replaced by $B_i$. In view of Proposition 3, this shows that
\begin{equation}\tag{8.8}
\left\{
\begin{aligned}
  \underline{T}_{1}^{1*}(A;B)&=\Extsh_{\Osh_X}^{*}(A,B),\\
  T_{1}^{1*}(A;B)&=\operatorname{Ext}_{\Osh_X}^{*}(X;A,B).
\end{aligned}
\right.
\end{equation}
In particular, taking $r=s=1$ and $A_1=\Osh_X$ in (8.6) recovers Theorem 3.
Formula (8.8) also implies
\[
  T_{0}^{1*}(B)=H^{*}(X,B),
  \qquad
  T_{1}^{0*}(A)=\operatorname{Ext}_{\Osh_X}^{*}(X;A,\Osh_X).
\]

\medskip
\noindent 3. Formula (8.1) shows that in general the functors
$\underline{T}_{r}^{s*}$ and $T_{r}^{s*}$ have components in both positive and
negative degrees. Using Remark 2, one sees that if $X$ has dimension $n$, then
the nonzero components of $\underline{T}_{r}^{s*}$ lie between $-(s-1)n$ and
$rn$ if $s>0$, and between $0$ and $rn$ if $s=0$; the nonzero components of
$T_{r}^{s*}$ lie between $-(s-1)n$ and $(r+1)n$ if $s>0$, and between $0$ and
$(r+1)n$ if $s=0$ (and, unless I am mistaken, if $r>0$, even between
$-(s-1)n$ and $rn$, respectively between $0$ and $rn$).

\section*{BIBLIOGRAPHY}

\noindent[1] Cartier (Pierre). --- Les groupes
$\operatorname{Ext}^{s}(A,B)$, Séminaire A. Grothendieck : Algèbre homologique,
t. 1, 1957, no. 3.

\noindent[2] Grothendieck (Alexandre). --- Sur quelques points d'algèbre
homologique, \emph{Tohoku Math. J.}, vol. 9, 1957, pp. 119--183.

\noindent[3] Serre (Jean-Pierre). --- Faisceaux algébriques cohérents,
\emph{Annals of Math.}, vol. 61, 1955, pp. 197--278.

\noindent[4] Serre (Jean-Pierre). --- Sur la dimension homologique des anneaux
et des modules noethériens, \emph{Proc. Intern. Symp. on Alg. Number Theory}
[1955, Tokyo and Nikko]. --- Tokyo, Science Council of Japan, 1956,
pp. 175--189.

\clearpage
